\documentclass[11pt]{article}
\usepackage[utf8]{inputenc}
\usepackage{amsmath,amssymb,amsthm}
\usepackage{hyperref}
\usepackage{listings}
\usepackage{xcolor}
\usepackage[margin=1in]{geometry}

\lstset{
    basicstyle=\ttfamily\small,
    breaklines=true,
    frame=single,
    backgroundcolor=\color{gray!5},
    numbers=none,
    xleftmargin=4pt,
    xrightmargin=4pt,
    aboveskip=8pt,
    belowskip=8pt
}

\newtheorem{conjecture}{Conjecture}

\title{Empirical Investigation of an Open Conjecture:\\
Every odd integer greater than 5 can be expressed as the sum of a prime and twic}
\author{Agentic NL$\to$Lean~4 Pipeline \\ \small Job \#10}
\date{\today}

\begin{document}
\maketitle

\begin{abstract}
This report documents the empirical investigation of an open mathematical conjecture
that could not be formally proved or disproved in Lean~4 with Mathlib.
Numerical experiments were conducted to gather evidence for or against the conjecture.
The empirical verdict is: \textbf{\textcolor{green!70!black}{Empirically Supported}}.
The conjecture remains formally open.
\end{abstract}

\section{Conjecture Statement}

\begin{conjecture}
\begin{lstlisting}
Every odd integer greater than 5 can be expressed as the sum of a prime and twice another prime.
\end{lstlisting}
\end{conjecture}

\section{Status}

\begin{center}
\fbox{\parbox{0.8\textwidth}{%
\textbf{Formal Status:} OPEN --- no Lean~4 proof or disproof was found.\\[4pt]
\textbf{Empirical Verdict:} \textcolor{green!70!black}{Empirically Supported}
}}
\end{center}

The pipeline attempted formal verification in Lean~4 with Mathlib but was unable to
produce a compiling proof or disproof. Empirical testing was then conducted to gather
numerical evidence.

\section{Basic Empirical Testing}

The following output was produced by the basic numerical experiment:

\begin{lstlisting}
======================================================================
=== EXPERIMENT PLAN ===
======================================================================

Conjecture (Lemoine / Levy, 1963):
    Every odd integer n > 5 can be written as n = p + 2q,
    where p and q are (not necessarily distinct) primes.

Strategy — test from multiple angles:

  (1) EXHAUSTIVE CHECK for all odd n in (5, N_MAX].
      Verify that at least one representation p + 2q exists for each n.
      If any n lacks a representation, report it as a counterexample.

  (2) COUNT the number of representations R(n) for each odd n > 5.
      Under the conjecture R(n) >= 1. We will study the growth of R(n)
      (it is conjectured to grow roughly like n / (ln n)^2 on average),
      which provides *redundancy evidence*: if R(n) is large and
      growing, random-model style heuristics strongly support the
      conjecture asymptotically.

  (3) MINIMUM-PRIME analysis. For each odd n, find the smallest prime p
      such that (n-p)/2 is prime. Track how this min-p grows with n.
      A bounded, slowly growing min-p suggests the conjecture is robust;
      a sudden jump would hint at possible failure.

  (4) RANDOM LARGE-n SAMPLING. Test 10,000+ large random odd integers
      up to ~10^7 (not full exhaustive) to look for counterexamples
      far beyond the dense regime.

  (5) EDGE / BOUNDARY CASES. Check n = 7, 9, 11, ... explicitly.
      Check powers-of-two-neighbors, Mersenne-like values, large twin-
      prime neighborhoods, etc.

  (6) STATISTICAL TEST. Compare empirical mean of R(n) with the
      heuristic prediction n / (2 (ln n)^2) (standard prime-tuples
      heuristic). Correlation should be strong and positive.

If every odd n>5 up to our bound has >=1 representation AND R(n) grows,
this is very strong empirical support (the conjecture has been verified
up to ~10^10 in the literature; we will re-verify a solid slice).


[1] Exhaustive check for odd n in (5, 200000]  ...
  Odd integers tested: 99,997
  Total representation pairs (p,q) counted: 94,873,744
  Counterexamples found (R(n)==0): 0
  ✓ Every odd n in (5, 200000] has at least one representation.
  min R(n) = 1, max R(n) = 4304
  mean R(n) = 948.77, median R(n) = 853.0
  Max 'smallest prime p' needed: 829 at n=189947

[2] Heuristic comparison: R(n) vs n / (2 (ln n)^2)
  Pearson correlation (binned) between R(n) and heuristic: 0.9999
  Ratio R(n)/heuristic : min=2.559, max=3.125, mean=2.829

[3] Smallest-prime-needed analysis
  n=  189947: need smallest p= 829, q=(n-p)/2=94559  (prime? True)
  n=   93167: need smallest p= 709, q=(n-p)/2=46229  (prime? True)
  n=  171961: need smallest p= 683, q=(n-p)/2=85639  (prime? True)
  n=   96563: need smallest p= 601, q=(n-p)/2=47981  (prime? True)
  n=   57887: need smallest p= 601, q=(n-p)/2=28643  (prime? True)
  n=   96527: need smallest p= 601, q=(n-p)/2=47963  (prime? True)
  n=   87007: need smallest p= 593, q=(n-p)/2=43207  (prime? True)
  n=   54863: need smallest p= 577, q=(n-p)/2=27143  (prime? True)
  n=  170099: need smallest p= 577, q=(n-p)/2=84761  (prime? True)
  n=  112589: need smallest p= 571, q=(n-p)/2=56009  (prime? True)

[4] Random large-n sampling up to 10^7
  Random large odd integers tested: 12,000 (range 7 .. 10,000,000)
  Failures (no representation found): 0
  Max smallest-prime p used in sample: 733
  Mean smallest-prime p used in sample: 56.17

[5] Edge / boundary cases
  n=         7  ->  reps (first 3): [(3, 2)]
  n=         9  ->  reps (first 3): [(3, 3), (5, 2)]
  n=        11  ->  reps (first 3): [(5, 3), (7, 2)]
  n=        13  ->  reps (first 3): [(3, 5), (7, 3)]
  n=        15  ->  reps (first 3): [(5, 5), (11, 2)]
  n=        17  ->  reps (first 3): [(3, 7), (7, 5), (11, 3)]
  n=        19  ->  reps (first 3): [(5, 7), (13, 3)]
  n=        21  ->  reps (first 3): [(7, 7), (11, 5), (17, 2)]
  n=        23  ->  reps (first 3): [(13, 5), (17, 3), (19, 2)]
  n=        25  ->  reps (first 3): [(3, 11), (11, 7
... [truncated]
\end{lstlisting}


\section{Experiment Code (Basic)}

\begin{lstlisting}[language=Python]
import matplotlib
matplotlib.use("Agg")
import matplotlib.pyplot as plt
import numpy as np
import random
import math
from sympy import isprime, primerange, nextprime
from collections import Counter
import time

t0 = time.time()

print("=" * 70)
print("=== EXPERIMENT PLAN ===")
print("=" * 70)
print("""
Conjecture (Lemoine / Levy, 1963):
    Every odd integer n > 5 can be written as n = p + 2q,
    where p and q are (not necessarily distinct) primes.

Strategy — test from multiple angles:

  (1) EXHAUSTIVE CHECK for all odd n in (5, N_MAX].
      Verify that at least one representation p + 2q exists for each n.
      If any n lacks a representation, report it as a counterexample.

  (2) COUNT the number of representations R(n) for each odd n > 5.
      Under the conjecture R(n) >= 1. We will study the growth of R(n)
      (it is conjectured to grow roughly like n / (ln n)^2 on average),
      which provides *redundancy evidence*: if R(n) is large and
      growing, random-model style heuristics strongly support the
      conjecture asymptotically.

  (3) MINIMUM-PRIME analysis. For each odd n, find the smallest prime p
      such that (n-p)/2 is prime. Track how this min-p grows with n.
      A bounded, slowly growing min-p suggests the conjecture is robust;
      a sudden jump would hint at possible failure.

  (4) RANDOM LARGE-n SAMPLING. Test 10,000+ large random odd integers
      up to ~10^7 (not full exhaustive) to look for counterexamples
      far beyond the dense regime.

  (5) EDGE / BOUNDARY CASES. Check n = 7, 9, 11, ... explicitly.
      Check powers-of-two-neighbors, Mersenne-like values, large twin-
      prime neighborhoods, etc.

  (6) STATISTICAL TEST. Compare empirical mean of R(n) with the
      heuristic prediction n / (2 (ln n)^2) (standard prime-tuples
      heuristic). Correlation should be strong and positive.

If every odd n>5 up to our bound has >=1 representation AND R(n) grows,
this is very strong empirical support (the conjecture has been verified
up to ~10^10 in the literature; we will re-verify a solid slice).
""")

# ---------------- (1) EXHAUSTIVE CHECK ----------------
N_MAX_EXH = 200_000
print(f"\n[1] Exhaustive check for odd n in (5, {N_MAX_EXH}]  ...")

# Sieve of Eratosthenes up to N_MAX_EXH
sieve = np.ones(N_MAX_EXH + 1, dtype=bool)
sieve[:2] = False
for i in range(2, int(math.isqrt(N_MAX_EXH)) + 1):
    if sieve[i]:
        sieve[i*i::i] = False
is_prime = sieve  # boolean array

# Primes <= N_MAX_EXH
primes = np.nonzero(is_prime)[0]

# For each odd n > 5, count representations n = p + 2q
# We iterate q over primes with 2q < n, then check if n - 2q is prime.
# Vectorize per-n via accumulation:
odd_ns = np.arange(7, N_MAX_EXH + 1, 2)
R = np.zeros_like(odd_ns, dtype=np.int64)
min_p = np.full_like(odd_ns, -1, dtype=np.int64)  # smallest p such that (n-p)/2 is prime

# For each prime q, mark all odd n = p + 2q where p is prime.
# Equivalent: for each prime q, for each prime p (odd p or p=2), n = p + 2q.
# Iterate over q; then n - 2q must be prime (=p).
# We'll do it per-n by checking primes q up to (n-2)/2.

# Faster: for each n, iterate primes p in increasing order, test (n-p)/2 is prime.
# Use a boolean lookup `is_prime` for (n-p)/2 which is <= (n-2)/2 <= N_MAX_EXH/2.

# Precompute small-prime list (all primes up to N_MAX_EXH-4):
primes_list = primes.tolist()

# To keep it fast & vectorized: for each prime p, mark all odd n >= p+4 such that
# (n-p)/2 is prime. That is, n = p + 2q for every prime q >= 2.
# We'll loop over primes p (a lot of them), and for each p loop over q? Too slow.
# Instead: loop over p (outer) and q (inner over primes list), set R[n]+=1 and min_p.
# Number of pairs ~ pi(N)^2 / 2 which is ~9600^2 = 9.2e7. Borderline.
# We'll use numpy: for each prime p, compute all n = p + 2*primes_array <= N_MAX_EXH.

primes_arr = primes.astype(np.int64)

# map from odd n to index in odd_ns: idx = (n-7)//2
for p in primes_arr:
    ns = p + 2 * primes_arr
    ns = ns[(ns <= N_MAX_EXH) & (ns >= 7) & (ns % 2 == 1)]
    idx = (ns - 7) // 2
    # Count each representation
    np.add.at(R, idx, 1)
    # Track minimum p
    # For entries where min_p == -1 or p < min_p, set min_p = p
    current = min_p[idx]
    mask = (current == -1) | (p < current)
    min_p[idx[mask]] = p

missing = odd_ns[R == 0]
print(f"  Odd integers tested: {len(odd_ns):,}")
print(f"  Total representation pairs (p,q) counted: {R.sum():,}")
print(f"  Counterexamples found (R(n)==0): {len(missing)}")
if len(missing) > 0:
    print(f"  First counterexamples: {missing[:20].tolist()}")
else:
    print(f"  ✓ Every odd n in (5, {N_MAX_EXH}] has at least one representation.")

print(f"  min R(n) = {R.min()}, max R(n) = {R.max()}")
print(f"  mean R(n) = {R.mean():.2f}, median R(n) = {np.median(R):.1f}")
print(f"  Max 'smallest prime p' needed: {min_p.max()} at n={odd_ns[np.argmax(min_p)]}")

# ---------------- (2) R(n) vs heuristic ----------------
print("\n[2] Heuristic comparison: R(n) vs n / (2 (ln n)^2)")
# ... [truncated]
\end{lstlisting}


\section{Conclusion}

The conjecture remains formally open. Numerical experiments \textbf{support} the conjecture --- no counterexamples were found across all tested parameter ranges.
Further investigation (both formal and empirical) is warranted.

\end{document}
